\documentclass[11pt]{article}
\usepackage[margin=1in]{geometry}
\usepackage{amsmath,amssymb}
\usepackage{enumitem}
\usepackage[T1]{fontenc}

\title{COMP2300 Set 7 Practice Exam}
\author{Topics: Pipelining, Hazards, Forwarding, Stalls, Branch Prediction, Superscalar}
\date{}

\begin{document}
\maketitle

\noindent\textbf{Time:} 120 minutes \hfill
\textbf{Total:} 100 marks

\vspace{0.5em}
\noindent\textbf{Instructions}
\begin{itemize}[leftmargin=1.5em]
  \item Answer in English.
  \item Part A is multiple choice.
  \item Part B contains pipeline calculations and hazard reasoning.
  \item Part C contains adapted past-exam style questions with source tags.
  \item Assume a five-stage ARM pipeline unless stated otherwise.
\end{itemize}

\section*{Part A: Multiple Choice (30 marks, 3 marks each)}
Choose exactly one option for each question.

\begin{enumerate}[label=\textbf{A\arabic*.}]
  \item The primary benefit of pipelining is:
  \begin{enumerate}[label=(\Alph*)]
    \item Lowering ISA complexity
    \item Improving throughput by overlapping instruction stages
    \item Eliminating branch instructions
    \item Guaranteeing CPI \(= 1\) on all workloads
  \end{enumerate}

  \item Which statement is true about latency in a pipelined CPU?
  \begin{enumerate}[label=(\Alph*)]
    \item It always decreases for each instruction
    \item It may stay similar or increase because of pipeline-register overhead
    \item It becomes zero after the pipeline fills
    \item It is equal to IPC
  \end{enumerate}

  \item A read-after-write hazard is a:
  \begin{enumerate}[label=(\Alph*)]
    \item Structural hazard
    \item Control hazard
    \item Data hazard
    \item Cache-coherence hazard
  \end{enumerate}

  \item Forwarding helps because it:
  \begin{enumerate}[label=(\Alph*)]
    \item Replaces the register file
    \item Sends a result directly from a later pipeline stage to an earlier consumer
    \item Eliminates all control hazards
    \item Prevents all load-use stalls
  \end{enumerate}

  \item A load-use hazard often still needs a stall because:
  \begin{enumerate}[label=(\Alph*)]
    \item The ALU cannot add two operands in one cycle
    \item The load result becomes available only after the memory stage
    \item The branch predictor is too slow
    \item The decode stage cannot read registers
  \end{enumerate}

  \item In the lecture design, a bubble is:
  \begin{enumerate}[label=(\Alph*)]
    \item A branch target buffer entry
    \item An empty operation inserted into the pipeline
    \item A renamed register
    \item The same thing as forwarding
  \end{enumerate}

  \item Which structure primarily stores branch target addresses?
  \begin{enumerate}[label=(\Alph*)]
    \item BHT
    \item BTB
    \item ROB
    \item IQ
  \end{enumerate}

  \item A one-bit predictor is weak on loop termination because it:
  \begin{enumerate}[label=(\Alph*)]
    \item Cannot distinguish branch PC values
    \item Has no branch target address
    \item Changes direction too easily after a single unusual outcome
    \item Requires two pipeline stages to update
  \end{enumerate}

  \item A Smith predictor uses:
  \begin{enumerate}[label=(\Alph*)]
    \item 2-bit saturating counters
    \item Floating-point scores
    \item A stack of branch PCs
    \item Only compile-time information
  \end{enumerate}

  \item A 2-way superscalar processor ideally can achieve:
  \begin{enumerate}[label=(\Alph*)]
    \item IPC up to 0.5
    \item IPC up to 1
    \item IPC up to 2
    \item IPC up to 5
  \end{enumerate}
\end{enumerate}

\section*{Part B: Computational Problems (50 marks)}
\begin{enumerate}[label=\textbf{B\arabic*.}]
  \item \textbf{Ideal pipeline timing (8 marks)}\\
  A five-stage single-issue pipeline executes 20 independent instructions with no hazards.
  \begin{enumerate}[label=(\alph*)]
    \item How many cycles are needed if the pipeline starts empty?
    \item What is the CPI for this run?
    \item What is the IPC for this run?
  \end{enumerate}

  \item \textbf{Forwarding or stall? (8 marks)}\\
  Consider the instruction sequence:
\begin{verbatim}
i1: ADD R1, R2, R3
i2: SUB R4, R1, #1
i3: LDR R5, [R6, #0]
i4: ADD R7, R5, R8
\end{verbatim}
  \begin{enumerate}[label=(\alph*)]
    \item Which pair can be handled by forwarding without any stall?
    \item Which pair is a load-use hazard that still requires a stall?
    \item Briefly explain why.
  \end{enumerate}

  \item \textbf{Load-use stall accounting (8 marks)}\\
  Assume a five-stage pipeline with forwarding and hazard detection. A program executes 40 instructions, and 6 of them cause exactly one load-use stall each.
  \begin{enumerate}[label=(\alph*)]
    \item How many total cycles are needed if the ideal no-hazard execution would have taken 44 cycles?
    \item What is the effective CPI?
  \end{enumerate}

  \item \textbf{Branch penalty comparison (8 marks)}\\
  A branch is taken 30 times in a program.
  \begin{enumerate}[label=(\alph*)]
    \item If the branch misprediction penalty is 4 cycles, how many wasted cycles occur for 30 mispredictions?
    \item If early branch resolution reduces the penalty to 2 cycles, how many cycles are saved over those same 30 mispredictions?
  \end{enumerate}

  \item \textbf{One-bit and two-bit prediction (8 marks)}\\
  Consider the branch outcome sequence:
  \[
    T,\ T,\ T,\ T,\ N
  \]
  \begin{enumerate}[label=(\alph*)]
    \item If a one-bit predictor starts in state \texttt{Taken}, how many predictions are correct?
    \item If a two-bit Smith predictor starts in state \texttt{Strongly Taken}, what final state does it end in after the sequence?
  \end{enumerate}

  \item \textbf{Superscalar IPC (10 marks)}\\
  A 2-way superscalar processor executes a 24-instruction program in 15 cycles.
  \begin{enumerate}[label=(\alph*)]
    \item Compute the achieved IPC.
    \item Compute the achieved CPI.
    \item Briefly explain why the achieved IPC may be lower than the ideal maximum.
  \end{enumerate}
\end{enumerate}

\section*{Part C: Past Exam Questions (Source-Tagged) (20 marks)}
\begin{enumerate}[label=\textbf{C\arabic*.}]
  \item \textbf{No-forwarding vs forwarding pipeline (Source: exam24r.pdf / exam24px1.pdf, Q7a, adapted) (7 marks)}\\
  A short ARM sequence runs on a five-stage pipeline.
  \begin{enumerate}[label=(\alph*)]
    \item Explain how the answer changes when the CPU has no forwarding/hazard detection.
    \item Explain how forwarding changes the hazard picture.
    \item Name one case where a stall is still required even with forwarding.
  \end{enumerate}

  \item \textbf{Smith predictor accuracy (Source: 2023\_Final\_Exam.pdf, Q2D, adapted) (6 marks)}\\
  Explain why a 2-bit Smith predictor usually performs better than a 1-bit predictor on loops that are mostly taken and then exit once.

  \item \textbf{Pipeline depth and register cost (Source: exam2025-main.pdf, Q7 / Q14, adapted) (7 marks)}\\
  Briefly explain the tradeoff between deeper pipelines and pipeline-register overhead. Why can a deeper pipeline improve clock frequency but still fail to improve overall performance proportionally?
\end{enumerate}

\end{document}
